\worksheettitle
  {Домашняя практика: ответы}
  {\modulemark{M07-H} Равные промежутки}

\begin{teacherbox}[Критерии]
  В каждом решении отдельно фиксируются число промежутков и действие с
  длиной. Прибавление единицы допустимо только при переходе от промежутков
  к числу объектов с включёнными концами.
\end{teacherbox}

\section{Ответы}

\begin{enumerate}
  \item a) $15-6=9$; б) $42-18=24$; в) $11$ промежутков.
  \item a) $(13-4)\cdot6=54$ мм;
    б) $(37-25)\cdot6=72$ мм.
  \item a) $12-7=5$, $35:5=7$ см;
    б) $48:6=8$, $20+8=28$.
  \item $45:5=9$ промежутков; $9+1=10$ фонарей.
  \item Причина: единица прибавлена к числу промежутков. Верно:
    $20-8=12$ промежутков.
\end{enumerate}

\begin{resultbox}[Маршрут]
  При системном смешении объектов и промежутков назначьте M07-R. При верной
  модели и вычислительной ошибке дайте короткую дополнительную практику.
  После устойчивого результата можно предложить M07\(\star\).
\end{resultbox}
